\documentclass[11pt]{article}
\usepackage[margin=1in]{geometry}
\usepackage{amsmath,amssymb,amsthm,mathtools}
\usepackage[hidelinks]{hyperref}
\newtheorem{theorem}{Theorem}
\newtheorem{corollary}[theorem]{Corollary}
\title{Generalized Hurwitz certificates for the coupled gap-three Wronskians}
\author{Internal exact algebraic verification; no priority claim}
\date{17 November 2026}
\begin{document}
\maketitle
Put $F_n(z,b)=W_{n,n+b-1}(z)$, where the gap-three Wronskian $W$
is defined in the \href{GAP_THREE_WRONSKIAN_RECURRENCE.tex}{Wronskian note}.
Its proved coupled recurrence is
\[
 F_0=1,\qquad
 F_{n+1}=(4n+b+2)\bigl[(4n+b+1+z)F_n(z,b+1)
                -z\partial_zF_n(z,b+1)\bigr]+z^2F_n(z,b).       \tag{1}
\]
Each $F_n$ is monic of degree $2n$ with strictly positive coefficients
for $n\ge1,b\ge0$. The parameter coupling in (1) is retained throughout.
We prove total nonnegativity of the whole infinite generalized Hurwitz
matrix at a finite block of indices and positivity of its initial
Euclidean pivots at every index. Total nonnegativity at all indices
remains open.

\section{Conventions and the external theorem}
For $P(z)=\sum_{j=0}^N c_jz^{N-j}$, with $c_j=0$ outside $[0,N]$,
use the one-based matrix
\[
 H_3(P)=(c_{3j-i})_{i,j\ge1}.
\]
Split $P=f_0+f_1+f_2$, where
$f_r(z)=\sum_{j\equiv r\ (3)}c_jz^{N-j}$.
In the nondegenerate step-three generalized Euclidean algorithm,
\[
 f_{i+3}=f_i-\frac{h_i}{h_{i+1}}zf_{i+1},\qquad
 \deg f_i=N-i,\qquad h_i=[z^{N-i}]f_i,\quad 0\le i\le N.      \tag{2}
\]
Strict positivity of all $h_0,\ldots,h_N$ implies total nonnegativity
of the infinite $H_3(P)$ by Holtz--Khrushchev--Kushel,
\emph{Generalized Hurwitz matrices, generalized Euclidean algorithm,
and forbidden sectors of the complex plane},
\href{https://arxiv.org/pdf/1506.07379v1}{arXiv:1506.07379v1},
Theorem 16 (also Theorem 31).
Their Theorem 4 then excludes zeros in the open sector
$|\arg z|<\pi/3$.
The source's matrix convention and polynomial division algorithm are
exactly those above. Its statements and algorithm have been inspected;
this use is an imported theorem, not an independent review of that paper.
The exact cached source is
\href{sources/holtz-khrushchev-kushel-2015-v1.pdf}{retained with the verification input}.

\section{Whole infinite matrices at finitely many indices}
\begin{theorem}\label{thm:block}
For every integer $2\le n\le16$ and every real $b\ge0$,
$H_3(F_n(\cdot,b))$ is totally nonnegative.
This is an exact symbolic computation over $\mathbb Q(b)$ followed
by the cited generalized Hurwitz theorem; no sampling of $b$ or
truncation of the matrix is used.
\end{theorem}
\begin{proof}
Apply (2) to the coefficients generated by (1). At each of the fifteen
specified indices the algorithm terminates at degree zero, producing
exactly $2n+1$ nonzero rows of degrees $2n,2n-1,\ldots,0$.
The replay verifies every row identity in $\mathbb Q(b)[z]$.
For every leading coefficient, cancellation gives
\[
 h_i(b)=p_i(b)/q_i(b),\quad p_i,q_i\in\mathbb Q[b],\qquad
 p_i(0),q_i(0)>0,
\]
and every coefficient of both polynomials is nonnegative.
These finite polynomial certificates imply $h_i(b)>0$ for every
$b\ge0$ at each specified index. Theorem 16 therefore applies to
all minors of the infinite matrix.

For completeness the deterministic exact input, elimination algorithm,
identity replay, and sign verification are retained in
\texttt{gap\_three\_hurwitz.py} and
\texttt{gap\_three\_hurwitz\_verify.py}; the coefficient-degree and sign
profiles are in \texttt{gap\_three\_hurwitz\_certificates.json}.
The input is generated from the explicit recurrence (1), rather than
from a numerical root finder or a stored success flag.
At $n=2$, the full leading-coefficient list is
\[
 1,\quad 2(b+4),\quad 3b^2+21b+38,\quad
 \frac{(b+3)(b+6)(3b^2+21b+38)}{2(b+4)},\quad
 (b+2)(b+3)(b+5)(b+6),
\]
which provides a small directly inspectable instance.
The checker additionally reconstructs the input from the independent
auxiliary-Wronskian formula at $2\le n\le5$ and rejects deliberately
corrupted row and sign controls.
\end{proof}

\begin{corollary}\label{cor:block}
For every $1\le n\le16$ and every $\beta>0$, the polynomial
\[
 T_{n,n+\beta-1,3}(t)=\sum_{k=0}^n\binom nk(n+\beta-1)_{3k}t^k
\]
is squarefree. Consequently the consecutive positive Gamma moments
at orders $n-1,n$ of $X(A+BX^3)$ cannot both vanish whenever
$2\le n\le16$. The ordinary factorial case is $\beta=1$.
\end{corollary}
\begin{proof}
The case $n=1$ is linear. For $n\ge2$, a common root of the three
initial arithmetic polynomials $f_0,f_1,f_2$ in (2) would propagate
through every row of (2), reaching the nonzero constant $f_{2n}$.
This is impossible. A common nonzero root of the cyclic sections of
$W$ would give such a common root, by taking a cube root and multiplying
the sections by their corresponding powers of that cube root.
The exact rank theorem in the
\href{HIGHER_GAP_ADJOINT_RANK.tex}{higher-gap adjoint note} therefore
proves squarefreeness. The consecutive-Gamma-moment bridge is proved
there as well. Pure monomials and zero coefficients are immediate
from their nonzero positive factorial/Gamma moments.
Thus this fixed-index simplicity consequence uses the exact degree
descent directly and does not need the imported generalized Hurwitz theorem.
\end{proof}

The sector mechanism remains useful for the unbounded target.
At any index for which $H_3(W)$ is totally nonnegative, HKK Theorem 4
excludes zeros in $|\arg z|<\pi/3$.
A hypothetical repeated root gives a nonzero three-point rotation
orbit of zeros of $W$. Either a point enters this open sector, or the
orbit meets both boundary rays $\arg z=\pm\pi/3$.
In the latter case the common cube is real, so the three cyclic
sections $W_0,W_1,W_2$ have a common real root, contrary to the proved
real-root exclusion in the Wronskian note.
Thus the open-sector theorem suffices; no closed-sector strengthening
is needed. This implication does not prove total nonnegativity at
any additional index.

\section{Initial Euclidean pivots at every index}
Write $C_j(n,b)=[z^{2n-j}]F_n(z,b)$, with zero coefficients outside
the polynomial's degree. Reading (1) backwards from the leading term
gives
\begin{align*}
 C_j(n+1,b)-C_j(n,b)
  ={}&(4n+b+2)C_{j-1}(n,b+1)\\
     &+(4n+b+2)(2n+b-1+j)C_{j-2}(n,b+1).                  \tag{3}
\end{align*}
Here $C_0=1$, $C_j(0,b)=0$ for $j>0$, and negative subscripts give zero.
For each fixed $j$, exact polynomial summation of (3) gives
$C_j\in\mathbb Q[n,b]$. The checker proves the difference identity
and initial value, which establish the formula for every integer $n$
by induction. It does not infer the formula from finitely many values.
For example, putting $v=2n+b$ gives
\[
 C_1=nv,\qquad C_2=\frac{n(n+1)}6(3v^2-3v+2n-2).
\]
Define the following explicit polynomials using $C_1,\ldots,C_{11}$
obtained from (3):
\begin{align*}
 D_3&=C_1C_3-C_4,&D_4&=C_2C_4-C_1C_5,\\
 D_5&=C_5D_3-C_2(C_1C_6-C_7),\\
 D_6&=D_4(C_1C_6-C_7)-D_3(C_2C_7-C_1C_8).
\end{align*}
A direct use of (2) gives
\[
 h_0=1,\quad h_1=C_1,\quad h_2=C_2,\quad h_3=D_3/C_1,
 \quad h_4=D_4/C_2,\quad h_5=D_5/D_3,
 \quad h_6=D_6/(C_1D_4).                                  \tag{4}
\]
The next two pivots are conveniently expressed using the special minors
\[
 \Delta_p=\det(C_{3j-i})_{\substack{k\le i\le k+r-1\\1\le j\le r}},
 \qquad r=\lceil p/2\rceil,\qquad
 k=\begin{cases}2,&p\text{ odd},\\1,&p\text{ even}.\end{cases}
\]
As always $C_j=0$ for $j<0$.
Direct determinant expansion gives
$\Delta_3=D_3$, $\Delta_4=D_4$, $\Delta_5=D_5$,
$\Delta_6=D_6/C_1$.
The next two steps of (2) give
\[
 h_7=\Delta_7/\Delta_5,\qquad h_8=\Delta_8/\Delta_6.           \tag{5}
\]
These identities are also instances of HKK Lemma 18, but here the
checker verifies them by explicit polynomial expansion from (2),
without importing that lemma as a further assumption.

\begin{theorem}\label{thm:initial}
For every integer $n\ge4$ and every $b\ge0$, all nine initial
leading coefficients $h_0,\ldots,h_8$ in (2) are strictly positive.
For every integer $n\ge3$ and every $b\ge0$, all seven initial
leading coefficients $h_0,\ldots,h_6$ in (2) are strictly positive.
For $n=2$, all five existing leading coefficients are strictly positive.
\end{theorem}
\begin{proof}
$C_1,C_2>0$ on the stated range. Substitute $n=x+2$ in $D_3,D_4$
and $n=x+3$ in $D_5,D_6$.
The resulting polynomials in $x,b$ have respectively
$35,70,145,287$ nonzero coefficients, all strictly positive,
and strictly positive constant coefficients.
Their complete rational coefficient lists are retained in the same
certificate file. The checker derives the polynomials from (3),
verifies the identities exactly, and verifies every coefficient's sign.
This is a finite exact polynomial certificate applying to all
$x,b\ge0$, hence to the claimed unbounded set of indices.
Furthermore, at $n=x+4$ the determinants $\Delta_7,\Delta_8$
have respectively $425,651$ nonzero rational coefficients, all
strictly positive, with constant coefficients
$196134634651648000$ and $713930070131998720000$.
Their full coefficient lists and exact determinant identities are
verified and retained with the earlier certificates.
Equations (4)--(5) prove positivity. At $n=2$ the displayed list above
proves the assertion directly.
\end{proof}

\section{A scalar polynomial target for all sizes}
For every integer $p\ge1$, set $r=\lceil p/2\rceil$ and
$k=2r-p+1\in\{1,2\}$, and form the same explicit determinant
$\Delta_p(n,b)=\det(C_{3j-i}(n,b))$.
Consider the stronger coefficient assertion
\[
 \Delta_p(x+r,b)\text{ has nonnegative rational coefficients in }x,b
 \text{ and a strictly positive constant coefficient}.        \tag{6}
\]
This remains a conjectural target for general $p$.
The replay verifies (6) exactly for $1\le p\le8$.
If (6) held for every $p$, then at any fixed integer $n\ge2$
it would give $\Delta_p(n,b)>0$ for every $1\le p\le2n$ and
$b\ge0$. These are exactly the special minors in HKK Theorem 15,
so the entire infinite $H_3(F_n)$ would be totally nonnegative.
The sector and boundary argument would then prove all-index gap-three
simplicity in the physical parameter range.

Reversing and transposing the determinant identifies it formally with
the dual Jacobi--Trudi determinant for the partition whose conjugate
is $(p,p-2,\ldots,p-2r+2)$, under the specialization $e_j=C_j$.
This determinantal identity is a possible route to a positive-network
or factorial-Schur interpretation; it does not prove positivity of
this specialization. The full arithmetic-progression parameter
structure in (1) is still needed.

Theorem~\ref{thm:initial} does not control later pivots, and
Theorem~\ref{thm:block} does not extend beyond its enumerated indices.
A uniform proof of positivity or nonvanishing of the later pivots,
a generalized Hurwitz factorization, or an appropriate positive network
would still be needed to settle all-index gap-three simplicity.
General lower exponents, all univariate binomial Strong Factorial,
general FC and unrestricted SFC remain open.
All computations here are internal exact replay; independent replay,
formal verification, external mathematical review, and priority claims
are not supplied.
\end{document}
